\chapter{模型评测}
\label{chap:nlp-evaluation}

一套语言模型在固定试题上答对得更多，可能来自参数学习，也可能来自更长的提示、更多次
采样或一个更强的选择器。若把这些改变同时发生的结果统称为“模型能力提升”，我们就
无法判断下一阶段该改数据、训练还是推断。本章承接前面各阶段，研究怎样把观察转成
有范围的证据：固定被测对象，准备独立问题，规定判定办法，再报告差异及其不确定性。

评测（evaluation）是按预先规定的材料与规则观察系统行为；它不等于给一个总分。
这件事有六个具体困难：分数由参数与完整运行过程共同产生，需要固定协议；评测题可能
已经出现在训练数据中；基准覆盖哪些能力、各占多少权重，会改变排序；开放式回答没有
唯一参考，只能靠人工或模型评判，而评判本身有偏差；推理、置信度与安全行为需要专门
的判据；两个模型的差异要有统计区间，并能归因到具体阶段。本章依次固定协议、准备
材料并检查污染、构造概率、任务结果与偏好三类判定、检查推理与安全，最后进行统计
比较与阶段诊断。我们沿用通知、数学求解与多轮交流的例子，同时区分基础概率模型、
交互模型和带检索或工具的完整应用；后续应用章再展开跨度、译文、证据记录和执行状态
的具体标准。

\section{评测对象与协议设定}
\label{sec:nlp-eval-protocol}

比较两个分数之前，先要说清楚分数属于谁、按什么要求计分、花了多少计算。这里有三个
困难：同一组参数接上不同的提示、解码或工具，得到的是不同的被测对象；“答得好”需要
拆成可以逐条判定的条件，而合并条件的方式会改变结论；采样次数、生成长度与失败重试
都会改变分数，却常常不出现在报告里。下面依次用协议记号固定被测对象，把目标行为拆成
判据，并规定预算与复现记录。

\subsection{模型、推断过程与应用系统}
\label{sec:nlp-eval-unit}

把同一组参数接到两种解码器上，回答准确率可能改变。因此，被测对象至少应说明参数版本、
词元接口、提示与对话模板、解码和停止规则。检索器、外部资料、工具、验证器参与回答时，
它们也属于产生结果的过程。只给模型名称不足以恢复一次比较。

可以把协议记作$\mathcal P=(\bm\theta,\mathcal D,\mathcal I,\mathcal B,\mathcal J)$：
$\bm\theta$是模型参数，$\mathcal D$是被测请求的分布或样本，$\mathcal I$是输入与推断
过程，$\mathcal B$是预算约束，$\mathcal J$是判定器。这个记号用来固定比较条件，并不
假定五者相互独立。研究训练改变时尽量固定后四项；研究推断改变时固定参数并按相同预算
比较。若评测完整应用，就报告组合结果，同时用组件替换定位收益来源。

\subsection{目标行为与可观察判据}
\label{sec:nlp-eval-criteria}

对通知提问“校外来宾何时截止报名”，正确日期、时区和对象范围可由材料直接检查。
若要求“写得亲切一些”，则存在多个合法表达，单一参考字符串不能覆盖全部答案。
我们需要把要求拆成可判定的条件：日期是否正确、是否遗漏截止时间、是否添加未经支持的
承诺，以及语言是否适合预定读者。条件可以分别报告，也可以预先规定怎样合并。

设一条请求有$K$个必要条件，条件判定为$j_k(x,y)\in\{0,1\}$。单条请求的严格合格指示为
$\prod_{k=1}^K j_k$，因为任何必要条件失败都会使该请求不合格；在题集上对这个指示取平均，
才得到严格合格率。单条请求的逐条件完成比例为$K^{-1}\sum_kj_k$，
再对题集平均，反映通常完成了多少条件。假设每个回答都遗漏同一项重要截止条件，
逐条件平均可能很高，严格合格率却为零。这不是指标矛盾，而是验收对象不同。
HELM以多场景、多指标报告语言模型行为，为避免单分数掩盖差异提供了具体参照
\scite{nlp-intro-liang2023helm}。

HELM把一个场景表示为任务、领域和语言的组合，适配器把不同模型接口转成同一评测请求，再由各项指标读取原始回答。其七类核心指标包括准确性、校准、稳健性、公平性、偏差、毒性和效率；这些指标不天然等权，也不都能用于每个场景。场景、提示和聚合是协议设计，不需要训练一个“HELM模型”；若某项指标调用毒性分类器，则应另外说明该分类器的参数及标注来源。原报告把模型在核心场景上的平均可比覆盖从此前的$17.9\%$提高到$96.0\%$，改善的是评测覆盖，而非模型正确率。HELM的用途是为模型选择提供同协议、多指标证据；其初始核心范围主要是英语，不能直接代表中文、多轮工具执行或所有安全行为。

\subsection{预算、版本与复现记录}
\label{sec:nlp-eval-budget}

推理专题已经表明，多采样能增加找到正确候选的机会。比较时应同时记录输入长度、允许
生成长度、候选数、验证次数、工具调用和实际消耗。相同词元上限不一定意味着相同时间或
费用，不同词元化下的“同样一千词元”也未必包含相同文字。报告应说明选择哪种预算，
以及为何它对应当前使用条件。

除了平均用量，还应记录超限、无结束符、解析失败和拒绝回答的处理规则。只对成功返回的
样本统计时延会遗漏超时；把格式失败悄悄重试到成功，会增加未披露的预算。保留模型与数据
版本、提示、随机种子和逐请求记录，可以重放具体失败。固定随机种子有助于复现，但单次
随机运行仍不足以说明结果对种子的稳定程度。

\section{评测材料与独立性}
\label{sec:nlp-eval-data}

协议固定之后，分数的含义还取决于题目本身。这里有三个挑战：训练与提示开发可能已经
接触过评测题，分数高可能只是记住了答案；一份题集只覆盖部分能力，总分又依赖分组权重；
总分不能说明模型是否真正使用了题中的关键条件。对应的做法依次是按来源划分并检查
污染，按能力与条件组织基准体系，以及构造只改变单个因素的对照题。

\subsection{划分单位与污染检查}
\label{sec:nlp-eval-contamination}

预训练、SFT、偏好反馈和提示开发都可能接触评测题。即使测试文件没有直接加入训练集，
同一篇通知的改写、答案解析和另一语言版本也可能已经出现。因此，划分应按源文档、事件、
模板或问题族执行，再检查精确重复与语义近邻。逐行随机划分往往不足以阻断共同来源。
评测题被训练数据覆盖的现象称为\term{数据污染}（data contamination）：模型可能凭记忆
而不是凭求解答对这些题，分数于是高估了它在新题上的表现。

污染检查按能否访问训练数据分成两类。能访问训练语料时，直接比对评测样本与语料的文本
重叠；只能调用模型时，从模型对样本的概率行为推断它是否见过。两类方法都只能限定
所检查的范围：没有发现重叠，只说明指定规则未找到，不能证明模型从未接触过题目。
因此还常用训练截止时间之后出现的新题、保密题、动态生成且人工核验的问题，以及结构
相同但内容改变的对照。新题也要保证难度和标注质量可比，否则分数下降既可能来自少
记忆，也可能来自题本身更难。

\subsubsection{n元重叠检测}
\label{sec:nlp-eval-ngram-overlap}

最直接的检查是问：评测样本中有多少连续的词串原样出现在训练语料里？连续词串越长，
偶然重合的可能越小，重合就越像是同一来源。设评测样本经规范化后为词序列
$u_1,\ldots,u_L$，它的全部$n$元组构成集合$G_n(x)$；训练语料的全部$n$元组构成索引
$\mathcal I_n$。样本的重叠比例为
\begin{equation}
 \rho_n(x)=\frac{|G_n(x)\cap\mathcal I_n|}{|G_n(x)|},
 \label{eq:nlp-eval-ngram-overlap}
\end{equation}
$\rho_n(x)\geq\tau$时把样本标为污染。规范化先统一大小写、删去标点和多余空白，使排版
差异不掩盖重合；$L<n$的短样本则直接检查整个样本是否出现。不同报告采用的$n$和$\tau$
并不相同：GPT-3按数据集取$n$，即该题集样本词数的第5百分位数，并限制在8到13之间，
样本只要有一个$n$元组与任一训练文档相同就标为“脏”样本，相当于$\tau$取最小的正值
\cite{brown2020}；PaLM取$n=8$，要求问题、提示或答案中至少$70\%$的8元组在预训练语料中
出现过\scite{chowdhery2023palm}。前者偏向多报，后者偏向只报大段重合。

实现上，训练语料可能有万亿词元，不可能逐对比较。常用做法是对每个$n$元组计算定长
哈希值，分片存入哈希集合或布隆过滤器，再流式扫描评测样本查询。布隆过滤器只会多报、
不会漏报，因而略微抬高污染比例。检测结果有两种用法：训练前从语料中删去与评测集重合的
文档，即\emph{去污染}；训练后把题集分成“干净”与“可能污染”两个子集，比较模型在两者上的
得分。若可能污染子集的分数明显更高，说明记忆可能参与了作答。

用一个教学示例看规则的作用。一道40词的数学题有$40-13+1=28$个13元组，其中3个出现在
训练语料，$\rho_{13}=3/28\approx0.11$。按“任一重合即标记”的规则它是污染样本，按
$\tau=0.7$的规则则不是。GPT-3在训练语料去污染时还忽略了在10篇以上训练文档中出现的
13元组，因为这些多是常见习语或法律套话\cite{brown2020}；同样，若这3个13元组恰好是“请写出解题过程并给出最终答案”之类的
通用指令，标记就是误报；比较前应先去掉题集共用的模板文字。$n$取小了，常见短语会造成
大量误报；取大了，只改动几个词的副本就能逃过检查。

n元重叠的优点是机制透明、可以逐条复核；缺点是需要训练语料访问权，并且只认字面重合。
改写题干、翻译成另一语言或替换题中数字后，题目的解法与答案结构不变，n元组却几乎全部
改变。补救的办法是用句向量检索语义近邻，再由人工或模型判断近邻是否为同一道题，代价是
检索与判断都引入新的误差。

\subsubsection{基于成员推断的检测}
\label{sec:nlp-eval-membership}

多数模型不公开训练数据，评测者只能调用模型、读取词元概率。这时可以借用隐私研究中的
\term{成员推断}（membership inference）：判断一段文本是否属于模型的训练集。共同的依据是
训练会提高模型对见过文本的概率，见过的样本因而得到更高的似然。成员推断与基于损失的
隐私攻击属于同一类方法，一般设定与攻击分类可参考\BookChapterRef{chap:trust-privacy}{卷三“隐私保护”章}。

直接用平均对数概率$-H$作判据有一个问题：常见、简单的文本即使没见过，平均概率也很高。
Shi等观察到，未见过的文本通常含有少数模型很难预测的“意外”词元，而训练中见过的文本
在这些位置上的概率会被抬高\scite{nlp-eval-shi2024mink}。Min-K\%概率方法因此只看最难
预测的那部分词元。设样本有$L$个词元，$\mathcal K_k(x)$是条件对数概率
$\log p_{\bm\theta}(x_t\mid x_{<t})$最小的$\lceil kL\rceil$个位置，判据为
\begin{equation}
 \operatorname{MinK}_k(x)
 =\frac{1}{|\mathcal K_k(x)|}\sum_{t\in\mathcal K_k(x)}
 \log p_{\bm\theta}(x_t\mid x_{<t}),
 \label{eq:nlp-eval-mink}
\end{equation}
$\operatorname{MinK}_k(x)>\varepsilon$时判为“可能见过”。$k=100\%$时它退化为平均对数概率；
$k$越小，判据越集中于低概率的尾部，原文的常用取值为$k=20\%$。判据不需要训练数据，
也不需要训练新模型，只需一次前向计算。

教学示例：一段10个词元的题干，各位置的对数概率为$-0.1$、$-0.3$、$-2.0$、$-0.2$、$-4.0$、$-0.5$、$-0.05$、$-3.0$、$-0.4$、$-0.6$。平均为$-1.12$，$k=20\%$时只取最低的两项$-4.0$和$-3.0$，
$\operatorname{MinK}_{0.2}=-3.5$。若模型在训练中见过这段题干，最难的两个位置可能被抬到
$-0.8$与$-0.6$；这时最低的两项变为$-2.0$与$-0.8$，判据升到$-1.4$，变化$2.1$，而平均值
只从$-1.12$变为$-0.56$，变化$0.56$。见过与否的差别集中在少数难预测的位置，平均会把它
稀释，这正是只看尾部的理由。

阈值$\varepsilon$的含义是：在已知未见的文本上，有多大比例会被误判为见过。通常取一批
确定不在训练集中的文本，例如训练截止日期之后才发表的同类文章，按指定的误报率（如
$5\%$）选$\varepsilon$；也可以在成员与非成员都已知的检测集上报告ROC曲线下面积，Shi等为此
构造了WikiMIA：以2023年以后才创建的维基百科事件页面为非成员，以较早的事件页面为
成员。判据的弱点同样来自它的依据：高频、
模板化的文本即使没见过也得到高值；若成员与非成员两组文本在主题或时间上本来就不同，
判据可能只是区分了分布，而不是区分了是否训练过。因此成员推断的结果是需要复核的线索，
不是污染的证明，更适合在整份题集上统计，而不是对单道题下结论。

表~\ref{tab:nlp-eval-contamination-methods}比较了几种污染检查。它们针对不同的访问条件
与泄露形式，常组合使用。

\begin{table}[htbp]
 \begin{tabularx}{\linewidth}{@{}p{16mm}p{24mm}X@{}}
  \toprule
  方法 & 需要什么 & 能检出与漏检的情形\\
  \midrule
  n元重叠 & 训练语料 & 字面复制；漏检改写、翻译与换数\\
  语义近邻 & 语料与编码器 & 改写与翻译；近邻判断有误差\\
  成员推断 & 词元概率 & 无语料时的线索；受文本常见程度干扰\\
  新题对照 & 截止后的新题 & 绕开记忆；须保证难度可比\\
  \bottomrule
 \end{tabularx}
 \caption{污染检查方法的访问条件与检出范围。前两种在训练端检查文本重合，后两种从模型行为或题目时间间接判断。}
 \label{tab:nlp-eval-contamination-methods}
\end{table}

\subsection{能力覆盖与基准体系}
\label{sec:nlp-eval-coverage}

一套中文通知问答不能代表数学、代码或其他语言的能力；一套选择题也不能自动代表多轮
办事。应从拟支持的请求范围出发，按语言、领域、长度、难度、信息可见性和交互轮次
分组，再说明每组样本如何取得。公开的\term{评测基准}（benchmark），即题目、判定程序与
计分规则都固定的公开题集，用于可重复比较；场景样本用于检查目标使用条件。两者可以
互补，却不能仅靠数量相加就消除覆盖缺口。

\begin{table*}[t]
 \begin{tabularx}{\linewidth}{@{}p{22mm}p{28mm}p{36mm}p{32mm}X@{}}
  \toprule
  基准 & 考察能力 & 题型与规模 & 判定与指标 & 已知问题\\
  \midrule
  MMLU & 57个学科的知识与理解 & 四选一，测试集14042题 & 选项准确率 & 头部模型接近饱和；部分题答案有误；提示格式影响分数\\
  GSM8K & 小学数学文字题的多步计算 & 测试集1319题，答案为数 & 提取最终数值后精确匹配 & 接近饱和；解答广泛公开，污染风险高\\
  MATH & 竞赛数学，7个科目、5级难度 & 测试集5000题，答案为表达式 & 规范化后等价匹配；常用500题子集 & 等价判定依赖规范化程序\\
  HumanEval & 由函数说明生成Python代码 & 164题，附单元测试 & 执行测试，pass@$k$ & 题量小；测试不足，补测后通过率下降\\
  AIME & 高中数学邀请赛 & 每年两套各15题，答案为0到999的整数 & 精确匹配，常对多次采样取平均 & 题量很少，区间宽；往年题已公开\\
  GPQA & 研究生水平的生物、物理、化学 & 四选一，448题，Diamond子集198题 & 选项准确率 & 题量小；专家亦难，标注一致性有限\\
  IFEval & 可验证的格式与内容指令 & 541条提示，25类约束 & 程序逐项核验，提示级与指令级 & 只查形式，不查语义\\
  Chatbot Arena & 开放对话中的用户偏好 & 用户自拟提示，两匿名模型作答 & 成对投票，Bradley--Terry评分 & 用户与提示分布有偏；受回答风格影响\\
  \bottomrule
 \end{tabularx}
 \caption{常用语言模型评测基准。规模指公开发布数据中的测试部分；各基准的判定方式决定它测量的对象，接近饱和或可能污染时，分数差异的意义随之减弱。}
 \label{tab:nlp-eval-benchmarks}
\end{table*}

分组还影响总分。设组$g$的准确率为$\widehat a_g$、权重为$\omega_g\geq0$且$\sum_g\omega_g=1$，总分为
$\sum_g\omega_g\widehat a_g$。按测试样本量加权反映当前题集组成，按目标请求频率加权
才对应指定使用分布；对各组等权又是另一种问题。设模型甲在数学组和对话组的准确率为
$0.8$和$0.5$，模型乙为$0.6$和$0.7$（教学示例）：两组等权时二者都是$0.65$；数学组占
$70\%$时甲为$0.71$、乙为$0.63$；对话组占$70\%$时甲为$0.59$、乙为$0.67$。权重改变排序，
因此应同时报告组内结果及权重依据，而不是只保留一个平均数。

表~\ref{tab:nlp-eval-benchmarks}列出常用的公开基准。它们按能力划分：知识与理解、数学、
代码、专家问答、指令约束和开放对话；判定方式从选项匹配、数值匹配、程序执行到人工
投票各不相同。

MMLU以57个学科的四选一题考察知识与理解\cite{nlp-eval-hendrycks2021mmlu}；GSM8K与MATH
分别是小学文字题和竞赛题，前者答案为数，后者要求把\code{\textbackslash boxed}中的表达式
规范化后判等\scite{nlp-intro-cobbe2021,nlp-eval-hendrycks2021math}。HumanEval以执行
单元测试判定代码，用第\ref{sec:nlp-eval-reasoning}小节的pass@$k$计分
\scite{nlp-intro-chen2021}。GPQA的题目由领域专家编写，并经过验证使非专家即使联网搜索
也难以答对\scite{nlp-eval-rein2024gpqa}。Chatbot Arena让用户对匿名模型投票，评分方法见
第\ref{sec:nlp-eval-arena}小节。

题量小的基准需要特别注意统计误差。把准确率看作$n$道独立题上的比例，标准误约为
$\sqrt{\hat a(1-\hat a)/n}$。GPQA Diamond有198题，$\hat a=0.6$时标准误约$0.035$，
$95\%$区间宽约$\pm7$个百分点；AIME一年30题，$\hat a=0.5$时标准误约$0.09$，区间宽约
$\pm18$个百分点。对每题多次采样再取平均，可以减小生成随机性带来的波动，却不能减小
“只有这30道题”带来的抽样误差。

基准的寿命也有限。准确率接近上限后，剩余错误中有不少来自题目本身的标注错误，模型
之间的差距不再反映能力；题目公开越久，进入训练语料的可能越大。常见的应对是换用
更难的新版本（例如十选一的扩展版选择题）、只用训练截止日期之后的新题（例如最近一年的
AIME），以及为代码题补充测试用例。每次替换都改变了题集，新旧分数不能直接比较。

\subsection{对照题与行为探针}
\label{sec:nlp-eval-probes}

为了判断模型是否使用了通知中的限制条件，可以只把“校外来宾”改为“全体参加者”，
观察答案是否相应改变。也可以保持句意而替换地名、日期格式或语序，检查不该改变的
判断是否保持。这类成对题控制了一个因素，比仅看题集总分更容易定位失败。

对照成立的前提是改动没有暗中改变其他验收条件。把时间移到过去可能同时改变可执行性，
把人物换成不同身份可能改变授权。行为探针说明模型在特定干预下怎样响应，不直接揭示
内部因果机制；若把输出变化解释为某个注意头或训练阶段的作用，还需要相应干预证据。

\section{概率、任务结果与偏好判定}
\label{sec:nlp-eval-measures}

题目准备好之后，要决定怎样判定回答。被评价的对象有三类：模型给参考文本的概率，
回答是否完成了任务，以及在多个回答之间人们更偏好哪一个。三类判定的计分单位不同，
本节分别给出公式和判定器，并说明每种代理指标与任务合格之间的距离。

\subsection{留出似然与概率质量}
\label{sec:nlp-eval-likelihood}

预训练阶段最直接的观察是模型给未参与更新的词元分配多少概率。第\ref{sec:nlp-classical-perplexity}小节
已经定义了交叉熵与\term{困惑度}；这里把它写成评测协议。固定非空的计分位置集合
$\mathcal M$，以自然对数计算
\begin{equation}
 H=-\frac{1}{|\mathcal M|}\sum_{(i,t)\in\mathcal M}
 \log p_{\bm\theta}(w_{it}\mid c_{it}),\qquad \operatorname{PPL}=\exp(H).
 \label{eq:nlp-eval-ppl}
\end{equation}
$c_{it}$是该位置实际可见的条件。数值越小，参考词元概率的几何平均$\exp(-H)$越高；这不等于算术平均概率更高。
直接比较时，词元化、上下文截断、可见条件、掩码与计分单位须保持相同约定。
模型结构可以不同，只要计算的是相同条件下同一参考事件的概率。

词元化不同会改变分母与预测事件。跨词元接口时，可以在同一原始文本、同一覆盖范围和
明确边界条件下报告每字节负对数似然，但仍需说明字节编码及哪些开销计入。掩码恢复模型
的某个去噪损失也不自动等于自回归似然，若使用似然界或估计器，应单列其定义。
概率质量可以帮助筛选预训练参数，却不能证明模型遵守指令、答案正确或使用证据可靠。

计算留出似然时，模型参数冻结，参考序列只提供计分目标及合法的先前条件，不反向更新模型。长文本可滑动窗口读取已有前缀，重叠区域用于条件但不重复计分；在互不相连的小块中逐块重启上下文，则测量的是另一个条件概率。对选择题计算选项似然，也须声明是否含选项标签、是否按长度归一化，以及同义选项怎样处理。长度归一化改变排序目标，“平均每词元概率高”不等于“完整选项概率高”。

双向掩码模型没有普通左到右联合似然时，可以逐位置遮蔽参考词元，再用其余位置预测该词元，形成伪对数似然
\begin{equation}
 \operatorname{PLL}(w_{1:T})
 =\sum_{t=1}^{T}\log p_{\bm\theta}(w_t\mid w_{1:t-1},w_{t+1:T}).
 \label{eq:nlp-eval-pseudo-likelihood}
\end{equation}
每一项的条件不同，这些条件概率的乘积不是归一化的序列分布。PLL不新增训练参数，但通常需要$T$次掩码前向。Salazar等用它对语音识别的100条候选转写重排，RoBERTa-large把LibriSpeech test-clean的词错误率从$7.26\%$降至$5.05\%$\scite{nlp-eval-salazar2020masked}。这展示了语言可接受性评分的用途；PLL与自回归PPL仍是两个不同的指标。

\subsection{答案核验与生成质量}
\label{sec:nlp-eval-answer}

数学题有可规范化的数值答案，代码有可执行测试，通知问答有来源字段。判定器先约定答案
提取与规范化规则，再计算正确率；提取失败应有独立计数，不能在看过模型输出后逐条放宽
规则。单元测试通过说明覆盖的输入输出条件成立，未覆盖路径和规格仍可能出错。

开放生成需要同时检查信息与表达。字符串重合、语义相似、事实蕴含和人工可读性评判测量
的对象不同：更相似的句子可能保留同一个错误日期，更自由的合法改写也可能得不到高重合
分。应把关键事实与硬约束作为独立判据，再用替代指标观察表达差异。各类生成指标的
机制及局限在文本生成与转换章展开。

对可验证约束，可以在每条提示中列出多项要求，并分别核验词数、格式、关键词等条件。
IFEval是这一思路的代表：以可检查的指令约束减少完全依赖开放式主观评分的情形
\scite{nlp-eval-ifeval2023}。它把提示拆成约束身份与参数，例如“回答恰含三段”或“不得出现指定词”，模型生成完整回答后，由确定程序逐项判定。原版本包含25类可验证指令与541条提示，一条提示可以组合多项约束。关键词、段数及语言识别等检查器须有明确的解析约定；中文的“词数”不能未经说明套用英语空格分词。提示级严格指标要求所有约束通过，指令级严格指标按每项统计，对应第\ref{sec:nlp-eval-criteria}小节的两种验收对象。

宽松指标先移除常见Markdown修饰或首尾附加语，再检查变换后的文本；它可减少格式误判，也可能使原本违反词数的回答变成合格。原论文中GPT-4的提示级严格合格率为$76.89\%$、指令级为$83.57\%$，同一模型的两种口径相差近7个百分点，说明验收单位本身就会改变分数。IFEval适用于输出约束保持和格式回归检查，不检查通知日期是否有来源，也不说明各约束在真实使用中的重要性。

\subsection{人工比较与模型评判}
\label{sec:nlp-eval-judges}

开放式回答没有唯一参考，只能比较。偏好对齐追求的是在指定判断标准下更常被选择，
而不是提高奖励器自己的分数。独立评价应另取提示、回答和评判者，隐藏模型身份，并平衡
候选展示次序。若$A$胜、$B$胜和平局数量分别为$W,L,T$，一种预先声明的胜率为
$(W+T/2)/(W+L+T)$；同时报告三类计数，使平局处理可审计。第\ref{sec:nlp-align-independent-feedback}小节
已用这一口径检查对齐前后的模型。评判模型不能是训练中用过的奖励模型，否则评测与
训练共用同一套误差。

两个问题由此产生。第一，候选模型有许多个，两两比较的票数怎样汇总成一张排行榜，
排名的不确定性有多大；第二，人工投票昂贵，换成大模型评判以后，怎样识别和修正评判
模型自己的偏差。下面分别处理。

\subsubsection{成对比较与竞技场评分}
\label{sec:nlp-eval-arena}

Chatbot Arena是一个公开的众包平台：用户输入自己的问题，平台随机选两个模型作答，隐藏
其名称；用户读完两份回答后投票（$A$好、$B$好、平局或都不好），投票后才显示模型名称
\scite{nlp-eval-chiang2024arena}。这样得到的是大量“某次对局中谁胜”的记录，每一对模型
被比较的次数各不相同，许多模型对之间甚至没有直接对局。我们需要一个模型，用所有对局
共同估计每个模型的强度。

第\ref{sec:nlp-align-bradley-terry}小节的Bradley--Terry模型正是这样的模型。它的结论是：
若每个对象有一个实数得分，胜负概率只由得分差决定，
$\mathbb P(i\succ j)=\sigma(s_i-s_j)$；得分只确定到一个公共常数。第\ref{chap:nlp-preference-alignment}章用它为同一提示下的
两份回答建立奖励模型，这里把对象从“回答”换成“模型”，得分$s_m$表示模型$m$在平台提示
分布上的总体强度。

国际象棋的\term{Elo等级分}（Elo rating）是这一模型的在线版本\scite{nlp-eval-elo1978}。
每局之后，按赛前分数计算期望得分并更新：
\begin{equation}
 E_i=\frac{1}{1+10^{(R_j-R_i)/400}},\qquad
 R_i\gets R_i+K(o-E_i),\quad R_j\gets R_j-K(o-E_i),
 \label{eq:nlp-eval-elo}
\end{equation}
其中$o\in\{0,\tfrac12,1\}$是$i$的实际得分，$K$是步长。令$s=R\ln10/400$，就有
$E_i=\sigma(s_i-s_j)$，即Bradley--Terry胜率；而$o-E_i$恰是单局对数似然对$s_i$的梯度。
Elo更新因此是对Bradley--Terry对数似然做逐局随机梯度上升：在$s$尺度上以步长$\eta$
更新$s_i\gets s_i+\eta(o-E_i)$，换回$R$尺度即$K=400\eta/\ln10$。

在线更新的结果依赖对局顺序。设$A,B$初始都是1000分，$K=32$，两局各胜一次。若先$A$胜后
$B$胜，最终$A=998.5$、$B=1001.5$；顺序反过来，结果恰好对调。同样的数据给出相反的排名，
是因为每次更新只用当前一局，越晚的对局权重越大。棋手的实力随时间变化，这种“近期加权”
有其道理；而竞技场中一个模型版本的参数固定不变，所有对局应同等对待。

因此Chatbot Arena改用全部对局的极大似然估计。设第$t$局在模型$i_t$与$j_t$之间进行，
$o_t$是$i_t$的得分，对数似然为
\begin{equation}
 \ell(\bm s)=\sum_{t=1}^{N}\Bigl[o_t\log\sigma(s_{i_t}-s_{j_t})
 +(1-o_t)\log\sigma(s_{j_t}-s_{i_t})\Bigr].
 \label{eq:nlp-eval-bt-likelihood}
\end{equation}
这是一个无截距的逻辑回归：每局的特征向量$\bm x_t=\bm e_{i_t}-\bm e_{j_t}\in\mathbb R^{M}$
在$i_t$处为$+1$、$j_t$处为$-1$，$M$为模型数。似然只依赖得分差，需固定一个参考模型的
得分为0，或加一个小的二次正则。$\ell$是凹函数；只要胜负记录满足一个连通条件，即无论把模型怎样分成两组，每组都至少
在一局中胜过或平过另一组，极大似然解就存在且唯一，并与对局顺序无关。若某组模型从未
输给组外模型，它们与组外的得分差可以无限拉大，似然没有有限的最大点，这时需要加正则。平台再用$1000+\frac{400}{\ln10}\hat s_m$换算回Elo尺度，便于
与旧榜单对照。

三模型的教学示例：$A$对$B$十局六胜，$B$对$C$十局六胜，$A$对$C$十局七胜。数值求解
式\eqref{eq:nlp-eval-bt-likelihood}，固定$s_C=0$，得$\hat s_B=0.417$、$\hat s_A=0.834$，
换算后约为1145、1072、1000分。拟合胜率为$0.603$、$0.603$、$0.697$，与观察的$0.6$、
$0.6$、$0.7$接近但不相等：单一得分要求$\operatorname{logit}$胜率满足可加性，
$\operatorname{logit}0.6+\operatorname{logit}0.6=0.81$，而$\operatorname{logit}0.7=0.85$，
极大似然在三组对局之间折中。

排名需要区间。对局并非来自一个可以重复的实验，最常用的做法是把对局视为独立样本，用
自助法估计得分的抽样分布，过程见算法~\ref{alg:nlp-eval-arena}。两个模型的区间明显
重叠时，它们的名次差别不应被解读为可靠的强弱关系。

\begin{algorithm}[竞技场评分：Bradley--Terry极大似然与自助区间]
  \label{alg:nlp-eval-arena}
  \begin{algorithmic}[1]
    \Require 对局记录$\{(i_t,j_t,o_t)\}_{t=1}^{N}$，$o_t\in\{0,\tfrac12,1\}$；参考模型$m_0$；重采样次数$R$
    \Ensure 各模型得分$\hat s_m$及$95\%$自助区间
    \Function{FitBT}{$\mathcal G$}
      \State 以$\bm x_t=\bm e_{i_t}-\bm e_{j_t}$为特征、$o_t$为标签，最大化式\eqref{eq:nlp-eval-bt-likelihood}
      \State 平移得分使$s_{m_0}=0$；\Return $\bm s$
    \EndFunction
    \State $\hat{\bm s}\gets$\Call{FitBT}{全部$N$条对局}
    \For{$r=1,\ldots,R$}
      \State 从$N$条对局中有放回地抽取$N$条，得到$\mathcal G^{(r)}$
      \State $\bm s^{(r)}\gets$\Call{FitBT}{$\mathcal G^{(r)}$}
    \EndFor
    \State 对每个模型$m$，取$\{s^{(r)}_m\}_{r=1}^{R}$的$2.5\%$与$97.5\%$分位数作为区间
    \State 按$1000+\frac{400}{\ln10}s$把得分与区间换算到Elo尺度
  \end{algorithmic}
\end{algorithm}

用户投票还受回答风格影响：更长、标题和列表更多的回答更容易胜出。\term{风格控制}
（style control）把风格差异作为协变量放进同一个逻辑回归：
\begin{equation}
 \mathbb P(i\succ j\mid x)
 =\sigma\bigl(s_i-s_j+\bm\gamma^{\top}\Delta\bm z_{ij}\bigr),\qquad
 (\Delta\bm z_{ij})_k=\frac{z_{i,k}-z_{j,k}}{z_{i,k}+z_{j,k}},
 \label{eq:nlp-eval-style-control}
\end{equation}
$z_{i,k}$是模型$i$在本局回答的第$k$项风格特征。Chatbot Arena团队2024年发布的风格控制
采用四项：回答词元数、Markdown标题数、列表项数与加粗次数。按两者之和归一化使差值成为
相对量：500与520个词元之差比20与40个词元之差小得多。$\bm\gamma$与$\bm s$一起估计。排序改用$\bm s$，含义是“两份回答风格
相同时的强度”。风格控制与第\ref{sec:nlp-align-reward-hacking}小节的长度控制同一思路；它也带来一个需要
声明的选择：若较长的回答确实更有帮助，控制长度会把这部分正当收益一并去掉。

竞技场评分的优点是覆盖真实用户的开放问题，不需要参考答案，且模型很难针对未知的用户
提问专门调优。它的局限来自同一来源：投票用户与提示分布偏向技术人群和英语，不代表
特定应用的用户；用户凭阅读判断，难以核对事实与代码正确性；同一用户多次投票、提供方
提交多个版本择优上榜，都会使对局不满足独立同分布的假设。

\subsubsection{模型评判的协议与偏差修正}
\label{sec:nlp-eval-judge-protocol}

人工投票每条都需要读者时间，开发中的频繁比较难以负担。\term{大模型评判}（LLM-as-a-judge）
用一个强语言模型读取问题与回答，输出分数或胜负\scite{nlp-intro-zheng2023judge}。
评判有三种形式。\emph{单回答打分}按评分细则给每个回答1到10分，不需要对照对象，
便于给许多模型一次性排序，但分数尺度不稳定，不同批次的8分未必同等。\emph{成对比较}
让评判模型在两份回答中选优或判平，判断更稳定，但$M$个模型需要$O(M^2)$次比较，常改为
每个模型只与固定基线比较。\emph{参考答案引导}先把参考答案交给评判模型，适用于数学和
推理题：原研究发现，评判模型可能被候选中的错误推导带偏，即使它自己独立求解时能答对。
多轮题还须把先前问答一并交给评判模型，否则它无法检查本轮是否保持前文要求。

MT-Bench用80个多轮问题覆盖写作、推理、数学、编程等八类任务，检查GPT-4评判与人工判断
的关系；在不计平局的口径下，GPT-4成对评判与专家多数意见的一致率约$85\%$，略高于专家
彼此之间的$81\%$\scite{nlp-intro-zheng2023judge}。一致率高说明评判模型可以替代多数人工比较，但
三类偏差仍需逐项处理。

\emph{位置偏差}指评判模型偏向放在某一位置的回答。修正办法是交换顺序各判一次，只有两次
判断一致时才计胜负，不一致时记平局。教学示例：100对回答中，两次都判$A$胜的40对、都判
$B$胜的30对，其余30对随位置改变判断。按交换规则，胜率为$(40+30/2)/100=55\%$；若只用
$A$在前的一次判断，而评判模型在那30对上都选了前者，胜率会被记成$70\%$。

\emph{长度偏差}指评判模型偏好更长的回答。第\ref{sec:nlp-align-reward-hacking}小节介绍了
长度控制胜率：拟合一个以长度差为协变量的胜率模型，再报告长度差为零时的预测胜率
\scite{nlp-align-dubois2024lc}。简化写法为
\begin{equation}
 \operatorname{logit}\mathbb P(m\succ b\mid x)
 =\theta_m-\theta_b+\phi\tanh\!\Bigl(\frac{\ell_m(x)-\ell_b(x)}{\sigma_\ell}\Bigr)+\psi_x,
 \label{eq:nlp-eval-lc-winrate}
\end{equation}
$\ell_m(x)$是模型$m$在提示$x$上的回答长度，$\sigma_\ell$是长度差的标准差，$\psi_x$吸收
提示难度。\term{长度控制胜率}（length-controlled win rate）取
$n^{-1}\sum_x\sigma(\theta_m-\theta_b+\psi_x)$，即把长度项置零后的平均胜率。
式\eqref{eq:nlp-eval-lc-winrate}与竞技场的风格控制式\eqref{eq:nlp-eval-style-control}
结构相同：两者都在胜率的对数几率上加一个风格差异项，再报告把该项置零后的结果；
区别在于前者的对手是固定基线，长度差用$\tanh$压缩，后者同时拟合全部模型，特征差按
两者之和归一化。

\emph{\term{自我偏好}}指评判模型偏向与自己同源的回答。原研究只在部分评判模型上观察到
这种倾向，数据量有限，尚不能确定其因果。处理办法是避免让候选模型兼任评判，或用多个
不同来源的评判模型投票取多数。

评判结论还须报告与人工判断的一致程度。原始一致率会被偶然一致抬高，常同时报告
Cohen's $\kappa=(p_o-p_e)/(1-p_e)$，$p_o$为观察一致率，$p_e$为按两方各自判断比例计算的
偶然一致率。例如$p_o=0.8$、$p_e=0.5$时$\kappa=0.6$。人工校准子集应与正式题集来自同一
分布；评判提示、评判模型版本与解码设置都属于判定器$\mathcal J$，变动后须重新校准。
若用这些标签再微调一个评分器，训练题、损失和测试划分须另行说明。模型评判适合开放
表达的批量比较和初步筛查，事实、算术与动作授权仍应交给相应的核验程序。

表~\ref{tab:nlp-eval-observation-contracts}把几种评测方法放回同一请求与输出集合。它们没有一个共同的“最高分最好”尺度；预测概率、约束合格与偏好票是不同的事件，即使候选预算相同，也不能相互换算。

\begin{table*}[htbp]
 \begin{tabularx}{\linewidth}{@{}p{25mm}p{34mm}p{39mm}X@{}}
  \toprule
  方法 & 实际输入 & 参数与判定来源 & 验收对象及缺口\\
  \midrule
  留出似然 & 参考词元与可见条件 & 冻结语言模型；计分无训练 & 参考概率；不验指令或事实\\
  PLL重排 & 固定候选、逐位置掩码 & 冻结掩码模型；权重开发集选 & 语言可接受性；非联合似然\\
  HELM & 统一场景与模型适配 & 预设协议；部分指标用分类器 & 多指标画像；不代表全部语言与任务\\
  IFEval & 约束及完整回答 & 确定检查规则，无模型更新 & 逐约束或整提示合格；不验内容真实性\\
  人工比较 & 盲化同题候选 & 标注者及判断标准 & 目标人群偏好；仍可分歧\\
  竞技场评分 & 用户自拟提示、两匿名回答 & 用户投票；Bradley--Terry拟合 & 平台用户的总体偏好；受风格影响\\
  模型评判 & 问题、回答及评分提示 & 现成模型或另训评分器 & 代理偏好；须校准位置、长度与自我偏好\\
  \bottomrule
 \end{tabularx}
 \caption{概率、规则、综合协议与偏好评价的观察对象。固定任务、输入、候选预算和判定信息后，才能解释每项证据支持的结论。}
 \label{tab:nlp-eval-observation-contracts}
\end{table*}

\section{推理、稳健性与安全行为}
\label{sec:nlp-eval-reasoning-safety}

前一节的判定针对单个回答。推理与安全还需要另外几类观察：多次采样中是否找得到正确
答案、模型对自己答案的把握是否可信，以及在风险请求与攻击输入下行为是否守住边界。
本节依次处理候选覆盖与轨迹、置信度与校准、有害响应与过度拒绝、对抗条件。

\subsection{答案能力与轨迹可靠性}
\label{sec:nlp-eval-reasoning}

推理轨迹长度增加后，最终答案可能更常正确，也可能只是更常尝试。一次生成的正确率、
多候选中至少存在一个正确答案的比例，以及选择器最终交付的正确率，应分开统计。
第一个观察单次生成，第二个观察候选覆盖，第三个还包含验证器的能力。

设一道题独立同分布采样$n$次，其中$c$次通过固定验证器。对$n\geq k$，常用的
pass@$k$估计为
\begin{equation}
 \widehat{\operatorname{pass@}k}
 =1-\frac{\binom{n-c}{k}}{\binom nk},
 \label{eq:nlp-eval-passk}
\end{equation}
其中$n-c<k$时分子为零。式子计算从$n$条记录不放回抽$k$条时至少有一条通过的概率；
在独立同分布采样假设下，它是$k$次生成候选覆盖的无偏估计\scite{nlp-intro-chen2021}。
直接用$1-(1-c/n)^k$代入会低估覆盖，因为它把不放回抽样当成了有放回抽样。
例如$n=5,c=2,k=2$时为$1-3/10=0.7$。它没有规定怎样从两条候选中选对答案，
因此不是实际选择器的交付准确率。相关采样或自适应搜索还会改变其统计解释。

答案正确也不足以证明每一步解释可靠。应另检查中间结论、引用、程序执行记录与反事实
改题后的响应；从外部观察到的书面轨迹，不应直接当作全部内部计算的忠实记录。
报告推理能力时，答案、过程可检验性和成本分别提供证据，缺一项就限制相应结论。

\subsection{置信度与校准}
\label{sec:nlp-eval-calibration}

模型给出答案时，最好还能说明自己有多大把握：把握低的回答可以拒答、转交人工或触发
额外验证。要让置信度可用，须回答三个问题：语言模型的置信度从哪里来；怎样检查它与
实际正确率是否相符；不相符时怎样修正。

置信度有三种来源。第一种是\emph{词元概率}：选择题取各选项字母的概率并在选项间归一化，
短答案取答案片段的长度归一化序列概率。它只需一次前向，但需要能读取概率，且同义答案
会把概率分散到不同字符串上。第二种是\emph{采样一致率}：对同一题采样$n$次，以出现最多的
答案的占比作为置信度，与第\ref{sec:nlp-execution-self-consistency}小节多数投票的计算相同；
它按含义而不是按字符串计数，代价是$n$倍的生成。第三种是\term{口头置信度}（verbalized confidence）：
在提示中要求模型在答案后写出“置信度：80\%”。它不需要读取概率，可用于只返回文本的接口，
但数值来自生成，常集中在少数整十值上。

置信度与正确率相符称为\term{校准}（calibration）：在所有报告置信度为$p$的回答中，正确的
比例应为$p$。卷三已经讨论了校准的一般定义、可靠性图、期望校准误差与温度缩放，并指出
校准与区分能力是两回事：对每题都报告总体正确率的模型也是校准的，却无法区分难易，
可参考\BookChapterRef{chap:trust-reliability}{卷三“可靠性”章}。下面把这些工具用到语言模型上。

\term{期望校准误差}（expected calibration error, ECE）先把置信度区间$[0,1]$切成$B$个箱，
记第$b$箱的回答集合为$B_b$，箱内平均置信度为$\operatorname{conf}(B_b)$、实际正确率为
$\operatorname{acc}(B_b)$，再按箱内样本占比加权平均二者之差：
\begin{equation}
 \operatorname{ECE}=\sum_{b=1}^{B}\frac{|B_b|}{n}
 \bigl|\operatorname{acc}(B_b)-\operatorname{conf}(B_b)\bigr|.
 \label{eq:nlp-eval-ece}
\end{equation}
把各箱的$(\operatorname{conf},\operatorname{acc})$画在图上，就是\term{可靠性图}（reliability diagram），
完全校准的点落在对角线上。教学示例：100道四选一题按置信度分入四个等宽箱，样本数为
20、30、30、20，平均置信度为$0.35$、$0.55$、$0.75$、$0.95$，实际正确率为$0.40$、$0.50$、
$0.60$、$0.80$。代入式\eqref{eq:nlp-eval-ece}，
$\operatorname{ECE}=0.2\times0.05+0.3\times0.05+0.3\times0.15+0.2\times0.15=0.10$。
总体平均置信度$0.65$高于正确率$0.57$，高置信度箱的差距最大，图~\ref{fig:nlp-eval-reliability}
显示这是典型的过度自信。

\begin{figure}[htbp]
  \centering
  % 第九卷 8.4.2 节：四选一题集上的可靠性图（教学示例，非实验数据）。
% 100道题按置信度分入四个等宽区间，各区间样本数为20、30、30、20；
% 区间平均置信度为0.35、0.55、0.75、0.95，实际正确率为0.40、0.50、0.60、0.80；
% 按式 eq:nlp-eval-ece 计算 ECE=0.2*0.05+0.3*0.05+0.3*0.15+0.2*0.15=0.10。
\begin{tikzpicture}[figure picture]
  \begin{axis}[
    width=\linewidth,
    height=58mm,
    xmin=0.2,xmax=1,
    ymin=0,ymax=1,
    xlabel={置信度},
    ylabel={实际正确率},
    xtick={0.2,0.4,0.6,0.8,1.0},
    ytick={0,0.2,0.4,0.6,0.8,1.0},
    grid=major,
    legend style={at={(0.5,1.03)},anchor=south,legend columns=3,
      /tikz/every even column/.append style={column sep=4mm}},
  ]
    \addplot[ybar interval,fill=FigureBlueFill,draw=FigureBlue,line width=.6pt]
      coordinates {(0.2,0.40) (0.4,0.50) (0.6,0.60) (0.8,0.80) (1.0,0.80)};
    \addlegendentry{区间正确率}
    \addplot[FigureStroke,densely dashed,line width=.8pt,domain=0.2:1,samples=2] {x};
    \addlegendentry{完全校准}
    \addplot[only marks,mark=*,mark size=1.6pt,FigureRed]
      coordinates {(0.35,0.40) (0.55,0.50) (0.75,0.60) (0.95,0.80)};
    \addlegendentry{（平均置信度，正确率）}
    \draw[figure link,FigureRed,line width=.9pt] (axis cs:0.35,0.35) -- (axis cs:0.35,0.40);
    \draw[figure link,FigureRed,line width=.9pt] (axis cs:0.55,0.50) -- (axis cs:0.55,0.55);
    \draw[figure link,FigureRed,line width=.9pt] (axis cs:0.75,0.60) -- (axis cs:0.75,0.75);
    \draw[figure link,FigureRed,line width=.9pt] (axis cs:0.95,0.80) -- (axis cs:0.95,0.95);
    \node[figure note,anchor=south] at (axis cs:0.3,0.03) {$n_b=20$};
    \node[figure note,anchor=south] at (axis cs:0.5,0.03) {$n_b=30$};
    \node[figure note,anchor=south] at (axis cs:0.7,0.03) {$n_b=30$};
    \node[figure note,anchor=south] at (axis cs:0.9,0.03) {$n_b=20$};
  \end{axis}
\end{tikzpicture}

  \caption{四选一题集上的可靠性图（教学示例）。100道题按置信度分入四个等宽箱，柱高为箱内实际正确率，圆点为（箱内平均置信度，正确率），红色竖线是它与对角线的差距。高置信度的两箱落在对角线下方，模型过度自信；按样本占比加权平均这些差距得到$\operatorname{ECE}=0.10$。}
  \label{fig:nlp-eval-reliability}
\end{figure}

ECE的数值依赖分箱。箱太少，同一箱内先高估、后低估的偏差会相互抵消；箱太多，每箱
样本很少，箱内正确率本身波动大。即使模型完全校准，有限样本下$|\operatorname{acc}-\operatorname{conf}|$
也几乎总为正，ECE估计偏大：100道题分成10箱时每箱约10道，箱内正确率的标准误可达$0.15$。
报告ECE时应写明箱数、等宽或等样本量分箱，以及每箱样本数；比较两个模型时用同一分箱。

置信度系统性偏高时，可以在不改变答案的前提下修正。\term{温度缩放}（temperature scaling）
把选项的对数几率$z_k$除以一个温度$T>0$再归一化，
\begin{equation}
 p_k=\frac{\exp(z_k/T)}{\sum_{j}\exp(z_j/T)},
 \label{eq:nlp-eval-temperature-scaling}
\end{equation}
$T$在一份独立的留出集上最小化负对数似然得到\scite{guo2017calibration}。所有对数几率按
同一比例缩放，最大的仍是最大，答案与准确率不变，只有概率的集中程度改变。两选项对数
几率为$(2,0)$时，$T=1$给出$0.88$，$T=2$给出$\sigma(1)=0.73$。这里的$T$与第\ref{sec:nlp-inference-temperature}小节采样时的
温度形式相同，作用却不同：采样温度改变生成的随机性，校准温度只改变报告的概率。
口头置信度没有对数几率可缩放，可以在留出集上拟合一个从报告值到正确率的单调映射，
例如按箱取经验正确率。

语言模型上有两个具体现象。其一，预训练模型在格式合适的选择题上常常校准得较好，
Kadavath等发现大模型对选项字母的概率与正确率相当接近，还可以让模型评估自己答案
“为真”的概率\scite{nlp-eval-kadavath2022know}。其二，偏好对齐会改变这一点。GPT-4技术
报告在MMLU子集上观察到，预训练模型的概率校准良好，经后训练后校准明显变差
\scite{openai2023gpt4}。偏好优化的目标奖励受偏好的回答，并不要求词元概率等于答对
频率。Kadavath等对RLHF策略模型也观察到明显的校准失调，而对所有评测统一取温度
$T=2.5$就大体修复了它，说明失调主要表现为概率过于集中。因此比较校准时须区分基座
与对齐模型，不能把基座的校准结论直接用于对齐模型。Tian等发现，对这类经人类反馈微调的模型，适当提示
得到的口头置信度往往比词元概率校准得更好\scite{nlp-eval-tian2023just}。

置信度的最终用途常是\emph{选择性预测}：只回答置信度不低于阈值$\tau$的题，其余拒答。
\emph{覆盖率}是回答的题目比例，\emph{选择准确率}是已回答题中的正确比例；$\tau$从低到高
扫描，得到覆盖率—准确率曲线。沿用上面的100道题：全部回答时覆盖率$100\%$、准确率
$57\%$；只回答置信度最高的两箱时覆盖率$50\%$、准确率$(18+16)/50=68\%$；只回答最高箱时
覆盖率$20\%$、准确率$80\%$。曲线是否上升取决于置信度能否把对错排开，即区分能力；
校准决定的是阈值能否直接读成“预期正确率”。两种性质应分别报告。
表~\ref{tab:nlp-eval-confidence-sources}比较了三种置信度来源。

\begin{table}[htbp]
 \begin{tabularx}{\linewidth}{@{}p{20mm}p{28mm}X@{}}
  \toprule
  来源 & 代价 & 主要问题\\
  \midrule
  词元概率 & 一次前向 & 须能读取概率；同义答案分散概率；对齐后常偏高\\
  采样一致率 & $n$次生成 & 按含义合并答案；$n$小时取值粗\\
  口头置信度 & 一次生成；纯文本 & 数值集中在整十；依赖提示\\
  \bottomrule
 \end{tabularx}
 \caption{语言模型置信度的三种来源。三者都须在留出集上用可靠性图与ECE检查，必要时用温度缩放或单调映射修正。}
 \label{tab:nlp-eval-confidence-sources}
\end{table}

\subsection{有害响应与过度拒绝}
\label{sec:nlp-eval-refusal}

安全对齐同时面临两种失败：对明确禁止的请求提供有效帮助，对合法请求作出不必要拒绝。
分别准备按既定政策标注的风险请求集$\mathcal H$和正常请求集$\mathcal S$，报告有害帮助
比例和正常帮助成功率。若只统计风险集上的拒绝比例，一个对所有请求都拒绝的模型也会
显得优秀；若只统计正常帮助，又会漏掉风险边界。

XSTest通过安全请求及容易触发误判的表述检查过度安全行为，提示我们拒绝本身也需要
与请求含义对应\scite{nlp-eval-xstest2024}。它构造含有风险词却语义合法的请求，并配对
真正不安全的请求来检查边界：公开套件含250条安全提示、200条不安全对照，安全提示包括
比喻、游戏语境、公开人物信息等类型。人工把回答分为完整拒绝、部分拒绝与遵从，再按
请求类型汇总。原论文中，Llama-2-70B-Chat在原系统提示下对安全提示的完整拒绝达$38.0\%$；
Mistral-7B-Instruct的完整拒绝原为$0.8\%$，加入安全系统提示后升到$9.6\%$，说明提示属于被测对象。低安全集拒绝率也须
配合不安全集的行为一起看。

对于允许解释但不允许执行的请求，还应区分回答边界与动作边界。评判要依据明确政策、
上下文和用途，关键词命中不能单独决定风险。政策改变后应更新判定规则并保留版本，
两个政策下的拒绝率不能直接解释为模型进退。

\subsection{对抗条件与输入边界}
\label{sec:nlp-eval-adversarial}

正常提问通过以后，还要问：攻击者能改变哪些输入，使模型偏离已有要求？越狱评测通常
研究攻击提示诱导被禁止响应；提示注入则常把指令放在检索文档、网页或工具返回中，试图
让模型把低信任材料当作有权修改任务的命令。两者的可控位置、可见信息和成功条件应
分别说明。仅在用户提示上测试，不能覆盖外部材料引起的越权。

一次攻击评测要固定攻击者可访问的信息、尝试次数、反馈及判定器。JailbreakBench用
标准化威胁模型、提示与评分规则支持可比较的越狱实验\scite{nlp-eval-jailbreak2024}。
它固定目标模型的系统提示、聊天模板、攻击者可改的用户输入及查询预算；攻击算法提出
提示、读取允许的反馈，再由判定器检查回答是否提供了所要求的有害帮助。JBB-Behaviors
包括十类共100种有害行为；按行为统计至少一次攻击成功，与按所有尝试统计成功率有不同
分母。攻击成功率只对声明的攻击集合和预算成立，更强攻击或新接口仍需重新检查。
稳健性与安全性的通用理论可参考\BookChapterRef{chap:trust-security}{卷三“对抗安全”章}，
本章把要求落实到语言输入、上下文边界和可观察输出上。

判定器本身是重要误差来源。原论文收集200组越狱提示与回答，由三位专家标注并取多数，
另加100条取自XSTest的良性样本，用来比较六种判定器。GPT-4与开放权重的Llama-3-70B同专家
标注的一致率都在$90\%$以上；Llama Guard的一致率只有$72\%$，漏报率约$61\%$；基于规则的
判定器一致率仅$56\%$。作者最终选用Llama-3-70B：它可以固定版本复现，误报率也较低，
不易把良性回答算成越狱成功。判定误差与预算限制都应进入结论：没判出攻击成功，不等于系统
没有漏洞。

提示注入的验收还要追踪信息权限。对同一通知应用，可以只替换低信任材料中的一段文字，保留用户目标与工具权限，观察是否发生未经授权的要求变更；攻击成功对象可能是泄露、错误引用或越权动作，不只是输出了风险词。低信任内容参与模型条件计算，不代表它有权改写系统任务。一个仅测试用户提示越狱的实验，不能支撑检索材料和工具返回均安全的结论。

\section{统计比较与阶段诊断}
\label{sec:nlp-eval-analysis}

前几节得到的是固定协议下的一组分数。要据此作决定，还剩三个问题：两个模型的分数差
有多少可能只是题目抽样造成的；差异应归于哪个训练阶段或推断改动；报告要包含什么，
读者才能判断结论能否用于自己的场景。本节依次用成对差值与自助区间估计不确定性，用
受控改变定位阶段贡献，再规定完整报告的内容。

\subsection{差异估计与不确定性}
\label{sec:nlp-eval-uncertainty}

两个模型在同一批$n$条请求上比较时，先为每条计算差值$d_i=s_A(x_i)-s_B(x_i)$，再求
$\widehat\Delta=n^{-1}\sum_i d_i$。成对差值保留了同一请求的难度关联，比把两组分数
当作不相关样本更符合实验设计。若请求按事件或文档成组，应以这些组为抽样单位，不能
把高度相关的多个问句当成独立重复。

可以按原抽样单位重采样，报告成对自助法区间；其解释依赖测试样本能代表目标总体，
也依赖聚类结构处理正确。仅有固定题集时，区间反映该抽样模型下的不确定性，不覆盖未来
分布变化、未计入的提示选择或未观测风险。多次比较后只挑最高分，也会引入选择偏差，
最终确认应使用独立材料。

\subsection{阶段贡献与失败归因}
\label{sec:nlp-eval-attribution}

比较预训练起点、SFT后、偏好对齐后与推理强化后的参数时，应使用共同可运行的协议，
并另报告各阶段推荐的输入模板。共同协议帮助定位参数变化，推荐模板观察实际用法，
二者不能相互替代。基座不遵守对话模板不必然表示其语言概率建模更差，聊天模型更受偏好
也不说明领域知识增加；第\ref{sec:nlp-eval-calibration}小节的校准变化就是对齐收益之外
需要单独记录的一项。

要判断某项数据或算法的作用，需要只改变该因素，或清楚声明哪些因素无法分开。
多个改动一起产生的收益只能归于组合。再按错误类型追踪：缺知识时考虑继续预训练，
格式与任务响应不稳定时考虑示范，评价标准不符时检查偏好反馈，候选覆盖足而选择错误时
检查验证器。这些是由证据支持的诊断方向，最终还需干预实验检验。

\subsection{完整报告与使用范围}
\label{sec:nlp-eval-report}

一份可使用的报告应让读者恢复被测对象、题集来源、信息边界、预算、判定、组内结果与
失败记录。样本数、分母、无效响应和区间与总分一起给出，才能解释“提高了多少”。
对重要约束还需列出绝对水平：假设错误率从$20\%$降到$10\%$，下降一半不表示剩余错误
已经满足一个要求错误率低于$1\%$的用途。这是构造算例，说明相对改进与可接受阈值必须
分别回答。

报告的最后应说明哪些语言、长度、时点、任务和攻击条件已有证据，哪些尚未覆盖。
模型进入新场景后继续收集失败，但新近用于调参的记录不能又当成独立测试；持续评测
需要定期补充新题并隔离旧题。评测由此成为训练与应用之间反复更新的证据接口，而不是
模型完成后永久有效的一张成绩单。

\section{本章小结}
\label{sec:nlp-eval-summary}

语言模型的评测结论同时依赖参数、材料、推断过程、预算和判定规则。评测材料先要检查
独立性：有训练语料时用n元重叠比对，没有时用Min-K\%概率一类的成员推断寻找线索，再以
新题对照；基准按能力组织，题量与饱和程度限制了它们能区分的差距。留出似然观察概率
建模，任务核验观察交付结果，偏好比较观察指定标准下的选择：竞技场评分用Bradley--Terry
极大似然汇总对局，并用自助区间与风格控制说明排名的不确定性和风格影响；模型评判通过
交换位置、长度控制和更换评判来源修正偏差，并以人工一致率校准。pass@$k$度量候选覆盖，
校准检查置信度能否当作正确率来读，安全与对抗测试观察行为边界。各类证据互补，不能
用其中一项替代其余对象。

前面各阶段现在具有了共同的比较依据：改变学习环节时固定可比较条件，改变推断方法时
把额外计算计入预算，再以分组结果、反例和不确定性限制结论。后续应用章将把同一原则
落实到文本结构、知识记录、生成内容、交互行动与语音时间条件，说明一般模型怎样才算
完成一个具体任务。

\paragraph{延展阅读}
HELM适合继续研究多场景、多指标的统一报告，IFEval把部分指令要求转成可检查规则，
Chatbot Arena给出成对偏好评分的平台与统计方法，XSTest与JailbreakBench分别考察过度
拒绝和越狱条件。对照阅读时，应先写清每项工作评价的行为对象、请求或攻击分布、判定器
和统计单位，再判断它能否支持当前用途
\scite{nlp-intro-liang2023helm,nlp-eval-ifeval2023,nlp-eval-chiang2024arena,nlp-eval-xstest2024,nlp-eval-jailbreak2024}。
